\section{Introduction}
\label{sec:intro}

\begin{figure}[t]
  \centering
  \includegraphics[width=\linewidth]{figs/teaser_3dps_v9b.pdf}
  \caption{\textbf{The right primitive for radar NVS.} Top: three method families on the three properties radar NVS demands. MC ray-tracing renderers (mmIR, Sionna RT) are physically faithful but too slow for multi-viewpoint optimization; optical-NVS adaptations (RadarSplat, Radar Fields, DART) are fast but discard phase. 3DPS is the first to satisfy all three. Bottom: 3DPS outputs on the held-out test frame of S1\,F185 --- $|\mathrm{RA}|$ render, raw ADC magnitude (inverse range FFT), and the optimized point cloud colored by learned $\varepsilon_r'$.}
  \label{fig:teaser}
\end{figure}

Millimeter-wave (mmWave) radar penetrates fog, rain, and airborne dust~\cite{guan2020through}, returns Doppler velocity directly, and is an essential sensor for robust perception in autonomous driving and robotics~\cite{gao2019experiments}. Its captures are also expensive: every change in transmitter--receiver geometry, chirp count, or carrier configuration forces full recapture, because each downstream radar product (range--azimuth, range--Doppler, raw ADC) is reconstructed from the same complex intermediate-frequency (IF) data via sensor-specific FFT chains. Novel view synthesis (NVS) for radar would address this directly --- a sparse set of real captures, an explicit and compressible scene representation, and a renderer that emits the \emph{complex} signal that downstream FFT pipelines consume (Fig.~\ref{fig:teaser}). No existing radar-NVS method does this.

Recent work approaches the problem from two sides, neither sufficient on its own. Optical-NVS adaptations~\cite{huang2024dart, 10.1145/3641519.3657510, takawale2025spinr, kung2025radarsplat} port primitives from RGB rendering --- neural radiance fields, hash-grid neural fields, and 3D Gaussians~\cite{mildenhall2021nerf, 3dgs} --- and supervise them on post-FFT power images, forcing the network to absorb antenna patterns, sidelobes, and FFT windowing artifacts into opaque learned features and discarding phase by construction. Physics-based differentiable renderers~\cite{sionna, hofmann2025inverse, chen2025invertwin, mmIR2026} model the radar forward model correctly but rely on Monte Carlo (MC) ray tracing, which makes the multi-viewpoint optimization that NVS demands intractable: the closest prior work, mmIR~\cite{mmIR2026}, requires ${\sim}$20 minutes per single viewpoint at 24k surface interactions per iteration. Neither family produces a renderer that is simultaneously physics-faithful, multi-viewpoint-tractable, and complex-valued.

We argue this is a primitive-mismatch problem. The radar forward model is angularly deterministic at every transmitter--receiver (TX--RX) pair: incidence and observation angles, antenna gains, and path lengths are computable in closed form, so the bidirectional scattering distribution function (BSDF) can be \emph{evaluated} rather than sampled, and the resulting phasor accumulates phase by integration rather than by learning. The right primitive is therefore an oriented 3D point carrying ITU-R P.2040 material parameters --- not an optical-radiance primitive emitting view-dependent RGB via spherical harmonics --- and the right rendering operation is splatting that phasor into range bins through a precomputed sensor point spread function (PSF). The 3DGS-era arguments against point primitives in optics --- ``holes, aliasing, strict discontinuity''~\cite{3dgs} --- do not bind at radar's $\Delta r{=}5.93$~cm range / ${\sim}1.4^\circ$ azimuth resolution, and the MC ray tracers we replace already compute scene hit points to synthesize ADC, so promoting those hits to persistent learnable primitives is a natural choice. We call this differentiable 3D Point Splatting (3DPS).

3DPS evaluates the BSDF and splats analytically, fused into two CUDA kernels with a third for the analytical backward pass. There is no MC sampling, no learned antenna pattern, no background subtraction, and no per-product retraining. Because each oriented point splats a complex phasor whose phase is integrated from geometric path length and material BSDF, the renderer produces complex range profiles directly --- not magnitude images --- and the same optimized scene therefore renders range--azimuth maps and raw ADC streams via standard FFT pipelines without retraining. On 20k oriented points across 8 training viewpoints, 500 Adam iterations complete in ${\sim}$3~minutes on a single RTX~4090: ${\sim}$50$\times$ faster than mmIR's per-scene MC ray tracing cost (${\sim}$160~min: $8\times{\sim}20$~min/viewpoint sequential), with all 8 training poses optimized jointly. On six outdoor ColoRadar~\cite{kramer2022coloradar} scenes with one held-out interpolated viewpoint per scene, 3DPS achieves \textbf{0.587} mean Pearson correlation on novel-view $|\mathrm{RA}|$ --- $1.8\times$ RadarSplat (0.333), $4.9\times$ Radar Fields (0.120), $7.2\times$ DART (0.081) --- with physically interpretable per-point materials.

Our contributions are:
\begin{enumerate}[leftmargin=*,itemsep=2pt]
\item \textbf{A primitive-mismatch argument for radar NVS.} We show that optical-radiance primitives (3D Gaussians, neural fields) are structurally unsuited to radar's deterministic angular forward model, and that an oriented 3D point carrying ITU-R P.2040 material parameters is the right replacement.
\item \textbf{A differentiable physically based point renderer that emits complex signals.} Closed-form BSDF evaluation and Hann-windowed PSF splatting in fused CUDA kernels; phase emerges from path-length integration rather than from learning. 20k oriented points across 8 training viewpoints, 500 Adam iterations in ${\sim}$3~minutes per scene on a single RTX~4090.
\item \textbf{Sparse-view radar NVS at the state of the art.} On six real outdoor ColoRadar scenes, \textbf{0.587} mean held-out novel-view $|\mathrm{RA}|$ Pearson correlation --- $1.8\times$/$4.9\times$/$7.2\times$ RadarSplat/Radar Fields/DART --- with physically interpretable per-point materials.
\end{enumerate}
